\documentclass[11pt]{article}
\usepackage[spanish]{babel}
\usepackage[margin=1in]{geometry}
\usepackage{amsmath, amssymb, amsfonts}
\usepackage{booktabs}
\usepackage{enumitem}
\usepackage{tcolorbox}
\usepackage{hyperref}
\usepackage{listings}
\usepackage{xcolor}
\usepackage{parskip}
\usepackage{pgfplots}
\usetikzlibrary{arrows.meta}

\pgfplotsset{compat=1.18}
\lstset{
    language=Python,
    basicstyle=\ttfamily\footnotesize,
    breaklines=true,
    columns=flexible,
    keepspaces=true,
    commentstyle=\color{gray}\itshape,
    keywordstyle=\color{blue!70!black}\bfseries,
    stringstyle=\color{green!40!black},
    frame=L,
    xleftmargin=2em,
    showstringspaces=false,
    tabsize=4
}

\title{\textbf{Laboratorio 5: Reinforcement Learning } \vspace{1em}\\
Inteligencia Artificial - CC3085}
\author{José Antonio Mérida Castejón}
\date{\today}
\begin{document}
\maketitle

\subsection*{Task 1 - Teoría}
\textit{Responda con criterio y análisis de ingenierx. No se esperan definiciones de libro, sino que analice las
consecuencias de las decisiones de diseño.}

\subsection*{El Dilema del Restaurante}

\textit{Considere lo que se habló en clase sobre Exploración y Explotación. Además considere que usted llega a una
ciudad nueva por 30 días. Hay 50 restaurantes. El primer día va al Restaurante A y la comida es excelente
(Reward = 10). Con esto en mente responda:}

\begin{itemize}

\item\textit{Si usted fuera un agente ``Greedy'' (Codicioso), ¿qué haría los 29 días restantes? ¿Por qué esta
    estrategia podría ser sub-óptima a largo plazo?}

    \noindent Si fuésemos un agente \textit{greedy}, la decisión sería de ir al \textit{Restaurante A} los 29 días restantes. Si ya visitamos
    un restaurante bastante bueno, no tendríamos razón alguna para arriesgarnos y potencialmente ir a un restaurante malo los días que restan. Esta
    estrategia puede llegar a ser subóptima, ya que realmente jamás exploramos la ciudad y es posible que encontremos mejores restaurantes. Si
    decidiéramos explorar más (y ser \textit{menos greedy}), puede que el restaurante vecino tenga una comida \textit{excepcional (Reward = 100)}.

\item\textit{Si implementa una estrategia $\epsilon-Greedy$ con $\epsilon = 0.1$, ¿cómo cambiaría su comportamiento? ¿Qué
    riesgo corre ahora?}

    En este caso, decidiríamos explorar los demás restaurantes en la ciudad. Con un valor de $\epsilon = 0.1$, 1 de cada diez veces dejaríamos de
    visitar el mejor restaurante hasta el momento para tomar una acción aleatoria (explorar algún otro). Aquí corremos el riesgo que el primer restaurante
    al que hayamos ido sea el mejor en general, dónde cada una de las veces que decidamos explorar vamos a estar tomando una decisión sub-óptima y 
    obteniendo rewards menores por explorar. En la práctica, es importante balancear los \textit{trade-offs} de exploración vs explotación cuando tenemos
    una situación dónde no tenemos visibilidad global del ambiente. Si tuviésemos esta información, el problema se vuelve trivial al visitar el restaurante
    con un mayor reward pero los entornos de la vida real no siempre cumplen esta propiedad.

\end{itemize}

\subsection*{Aprendizaje On-Policy vs Off-Policy}

\textit{Considere lo dicho sobre las formas de aprendizaje on vs off policy que se habló en clase. Considere con
esto el escenario donde un robot está aprendiendo a caminar, con cada una de estas formas:}

\begin{itemize}
\item\textit{Caso SARSA (On-Policy): El robot aprende basándose en los movimientos torpes que realmente hace
    mientras explora (incluyendo cuando se cae).}

\item\textit{Caso Q-Learning (Off-Policy): El robot aprende asumiendo que en el futuro siempre tomará la mejor
    decisión posible (aunque ahora se esté cayendo por explorar)}
\end{itemize}

\noindent\textit{Responda:}

\begin{itemize}
\item\textit{Si el entorno es muy peligroso (ej: un robot caminando al borde de un precipicio), ¿cuál algoritmo es más ``seguro'' durante el
    entrenamiento y por qué? (Hint: Piense en quién penaliza más la exploración riesgosa).}
\end{itemize}

El algoritmo \textit{SARSA} sería la elección óptima en este caso, esto se debe a que este algoritmo evalúa la poliza que está activamente siguiendo.
Por otro lado, \textit{Q-Learning} asume que continúa a lo largo de la ruta óptima.
\subsection*{Bootstrapping}
\textit{En clase mencionamos el término ``Bootstrapping'', repsonda:}

\begin{itemize}
    \item\textit{Explique en sus propias palabras (sin fórmulas) qué significa que un algoritmo actualice sus estimaciones basándose
        en otras estimaciones en lugar de esperar al final del episodio (como monte carlo). ¿por qué esto es vital para tareas continuas
        que nunca terminan (como el control de temperatura de un servidor)?}
\end{itemize}

Explicado mediante una analogía, el proceso de \textit{bootstrapping} es como corregir el rumbo mientras se va manejando en lugar de esperar
un choque para darse cuenta que algo salió mal. Esto es clave para tareas continuas, dónde no tenemos un ``punto final'' dónde podamos finalizar
y tener una retrospectiva del \textit{performance} a lo largo del task. En el ejemplo del servidor, la tarea realmente nunca finaliza y si
fuésemos a actualizar al finalizar la tarea estaríamos quemando múltiples servidores antes de corregir el comportamiento riesgoso.

\end{document}
